\section{Discussion}\label{sec:discussion}

The model isolates one mechanism: endogenous repricing by a shared private intermediary can change the local strategic relation between decentralized public investments. Several boundaries matter for interpreting that result.

First, the analytic result is local rather than global. Theorem \ref{thm:reversal} establishes the sign reversal for sufficiently small positive $\beta$ on a regular beta-zero B3 stationary branch and a regular G branch through the symmetric regular beta-zero central-interior full-game stationary state covered by the negative-cross-effect formula. The local argument also maintains support-side interiority, strict routing inequalities, nonsingular systems, the matched-price path, and strict private and public second-order conditions. It does not provide a primitive formula for the endpoint $\varepsilon$, prove an asymmetric beta-zero G result, prove global public optimality along every continuing branch, prove that such branches exist for every primitive vector, or characterize all primitives for which reversal occurs.

Second, the one-vector numerical exercise is deliberately weaker. Stage 11R corrects support-side surplus whenever participation saturates and re-runs the all-regime search with private repricing after G deviations, a fixed matched fee during B3 deviations, boundaries, low-investment histories, saturation neighborhoods, and multiple participation starts. No profitable public deviation is detected under that documented search. But the procedure does not establish a rigorous upper bound on regret over unsearched points or enumerate every possible downstream equilibrium. It is therefore search evidence, not a certified global-equilibrium or uniqueness result.

Third, the earlier draft vector is not global-equilibrium authority. It admits profitable finite deviations once routing regimes omitted by the original central-regime solver are restored. The exact rational Krawczyk calculation for that vector remains useful only as a local stationary-root diagnostic. It verifies local algebraic geometry, not global best-response optimality. The current vector is used only for the corrected numerical search and local welfare illustration.

Fourth, the novelty is proposition-level and exposed to a natural objection. Positive cross-side network feedback is familiar, negative effects transmitted through an endogenous follower price are familiar, public investments can already be strategic substitutes in fiscal-competition models, and third-party feedback can reverse strategic complementarity in other settings. The economic content here is narrower: with the private price level matched across the two environments, activating the shared private intermediary's optimizing price response can reverse the sign of the same public--public local best-response slope. We do not claim a general theory of feedback-induced strategic reversal.

Fifth, the repaired numerical vector uses a strong remote-public access friction satisfying
\[
\kappa_L+\tau>v+\alpha.
\]
Under this sufficient condition a nonresident rival public hub is dominated by non-participation for every project type. The condition prevents the direct rival-hub free-riding deviation that invalidated the draft vector. It also limits interpretation: the numerical exercise is not evidence of direct user poaching between public hubs. Public--public interaction at the reported states is mediated through the shared private route and cross-side participation. The numerical value of $\tau$ is constructive rather than empirically calibrated.

Sixth, primary-route single-homing is a modeling abstraction. Innovation projects and the firms behind them can interact with several networks. The maintained interpretation is narrower: a focal project has one principal intermediation route for the activity being modeled. Allowing project multihoming could attenuate route substitution and lies outside the present analysis.

Seventh, the private objective and instrument set are reduced-form approximations. Real commercial innovation communities may value occupancy, membership revenue, sponsorship, ecosystem growth, or longer-run strategic goals, and may use richer contracts than a single posted price. Similarly, a posted membership or access fee may bundle workspace and community services. The model uses a one-dimensional effective access price to represent the margin through which a commercial alternative responds to project demand; it does not claim robustness to arbitrary contracting environments. The sequence in which public facilitation is chosen before private pricing is likewise a maintained feature of the research question rather than a claim about all institutional timings.

Eighth, the institutional evidence establishes an analogue class, not an observed treatment-control relationship. We do not observe a causal episode in which expansion of a regional public hub induced CIC Tokyo, QWS, Level39, or another named operator to cut its price. The negative price response is a theoretical prediction. A useful empirical design would combine changes in regional facilitation capacity with project routing, partner participation, and flexible private access prices.

Finally, strategic substitution is not a policy diagnosis. The fixed-price environment is an identification benchmark, not a recommendation to regulate private prices. At the reported G stationary candidate, the corrected result pipeline finds a positive local national-welfare derivative in a coordinated $+x_i$ direction, with the own derivative computed numerically and the national derivative verified both directly and by decomposition. This is not a solved first-best problem, a global social optimum, or a general welfare theorem. No specific merger, subsidy, mandate, or price restriction is shown to be optimal.
